% ============================================================================
% 第 4 章 · 系统实现
% 目标：~10000 中文字 / 真 Rust 代码（贴自 fuxi 主仓）/ A2A 节占重头
% ============================================================================

\section{系统实现}

第三章给出了伏羲（fuxi）的整体架构与各模块对外契约。本章把视角下沉到代码层，结合主仓 \texttt{/Users/e0\_7/fuxi}（HEAD: \texttt{04eb34c}）的真实实现，依次说明开发环境、Rust workspace 组织、事件数据结构与 Lamport 因果序、任务派发主循环、A2A 协议从零实现、跨节点分布式通讯六个方面。代码片段均直接节选自 \texttt{crates/} 下相应文件，为版面做了必要的行宽折断与无关路径省略，但不修改控制流；行号与文件路径在每段代码上方注明，便于回到源码核对。本章的论述以工程取舍为主，不复述协议规范本身。

\subsection{开发环境与技术选型}

伏羲选择 Rust 作为唯一实现语言。这一选择基于三方面权衡：第一，Rust 兼具 C 级别的执行性能与编译期内存安全，避免了 Go/Java 等语言 GC 暂停带来的尾延迟尖峰，对于一个长时间常驻、需要为多个本地子进程做实时事件 republish 的编排平台，确定性的延迟分布比平均吞吐更重要；第二，Rust 的所有权与生命周期模型在编译期把"跨线程共享可变状态"这一典型并发陷阱明确化，使得事件总线、门客注册表（\texttt{Shelf}）等多生产多消费数据结构在协议层面就具备形式化保障，单元测试只需覆盖业务路径而无需重复验证内存安全；第三，Rust 异步生态在 2024 年已经成熟，Tokio 提供了工业级的多线程运行时与并发原语\cite{tokio2024broadcast}，Axum 提供了组合式的 HTTP/WebSocket/SSE 路由抽象，Reqwest 处理连接池与超时，sqlx 提供编译期 SQL 校验，这些库的稳定性使得伏羲可以专注于协议语义而不必重造底层基础设施。

异步运行时选用 Tokio 多线程版本（\texttt{tokio = \{ version = "1", features = ["full"] \}}）。由于伏羲内部同时承载事件订阅、HTTP 处理、子进程 stdio 读写、WebSocket 反连、cron tick 等多种并发模式，多线程运行时可以让 axum handler 与 dispatch pump、scheduler keeper 共享同一个 work-stealing 调度器。HTTP 服务统一选用 Axum 0.7 系列，因其 \texttt{Router::with\_state(...)} 模式无需引入依赖注入框架即可在 handler 间共享句柄；WebSocket 与 Server-Sent Event 由 Axum 原生提供\cite{fette2011websocket}，避免了引入第二个 HTTP 框架。持久化层选用 SQLite（通过 sqlx 异步访问），开启 WAL 模式以允许并发读写\cite{sqlite2024wal}，单一文件存储与事件溯源（event sourcing）的"append-only 日志即真相源"模型天然契合，且部署不需要任何外部数据库进程。前端 IM 端采用 PWA 形态，通过 Service Worker 与 manifest 提供桌面安装能力，在不依赖任何应用商店的前提下覆盖移动端与桌面浏览器。

值得特别说明的是 Axum 与 Tokio 的紧耦合带来的一个隐性好处：整个 fuxi-cli daemon 进程内只有一个 tokio runtime，所有子模块（fuxi-im 的 HTTP handler、fuxi-firehose 的 WebSocket handler、fuxi-scheduler 的 cron keeper、fuxi-orchestrator 的 dispatch pump）都跑在同一个调度器上，共享 worker pool。这避免了"多个 runtime 争抢线程"的常见陷阱（典型表现是某个 runtime 的 task 莫名其妙变慢），也让事件总线的 broadcast channel 跨模块零拷贝传递成为可能。要在 Rust 里达到这种"全进程一个 runtime"的状态需要团队级别的纪律——任何引入新依赖时必须确认它不会自己起一个 runtime（一些老 crate 会偷偷调 \texttt{tokio::runtime::Runtime::new}），fuxi 在依赖审查时把这条作为硬指标，引入任何新 crate 前都会用 \texttt{cargo tree -i tokio} 检查。

代码组织方面，伏羲采用单一 Cargo workspace，根目录 \texttt{Cargo.toml} 通过 \texttt{[workspace.dependencies]} 锁定所有共享依赖版本。这种 monorepo 化的依赖管理确保 13 个内部 crate 间的 \texttt{tokio}、\texttt{serde}、\texttt{thiserror} 等基础库版本完全一致，避免了符号冲突与重复编译。代码风格由 \texttt{cargo fmt} 强制统一，质量门禁为 \texttt{cargo clippy --workspace --all-targets -- -D warnings}，任何 lint 警告都会让 CI 红灯。完整测试集通过 \texttt{cargo test --workspace} 触发，截至本文写作时全 workspace 共有约 380 个单元测试与集成测试通过。

工具链版本通过 \texttt{rust-toolchain.toml} 固定到 stable 通道（\texttt{edition = "2024"}），避免不同开发者本地工具链漂移。错误处理统一使用 \texttt{thiserror} 派生 \texttt{Error} enum，library crate 内严禁 \texttt{unwrap()} 与 \texttt{panic!}，只允许在二进制 crate 的顶层错误边界（\texttt{main} 函数与 axum 中间件最外层）做最后兜底——这是为了避免一条 panic 让整个 daemon 崩溃带走所有门客状态。日志使用 \texttt{tracing} crate 的结构化 span，关键路径上 \texttt{\#[tracing::instrument]} 自动注入 span，便于在事故复盘时通过 trace id 还原跨进程调用链。整套工具链与代码风格的一致性确保了 13 个 crate 协同演进时的可维护性，新的开发者克隆仓库后只要通过 \texttt{cargo check} 即可在 5 分钟内验证本地环境完整。

\subsection{Rust Workspace 结构}

伏羲 workspace 共包含 13 个 crate，总代码规模 82\,515 行（统计自 \texttt{find crates -name "*.rs" \textbar{} xargs wc -l}，仅计 Rust 源文件，含集成测试与单元测试）。各 crate 按职责分层组织，自底向上依次是核心类型层、基础设施层、协议层、编排层、适配器层与集成层。表~\ref{tab:crate-layout} 列出了各 crate 的代码规模与主要职责。

\begin{table}[H]
\centering
\caption{伏羲 workspace 各 crate 代码规模与职责}
\label{tab:crate-layout}
\begin{tabular}{lrl}
\toprule
\textbf{Crate} & \textbf{LOC} & \textbf{主要职责} \\
\midrule
fuxi-cli & 33\,408 & 二进制入口 + REPL + daemon + dist controller \\
fuxi-im & 16\,149 & IM 后端 axum 服务（任务/对话/intervene 端点） \\
fuxi-orchestrator & 10\,220 & 玄女编排 trait + dispatch pump + 门客注册表 \\
fuxi-agent-cc & 3\,900 & Claude Code 门客适配器（WebSocket 反连） \\
fuxi-workspace & 3\,043 & git worktree 隔离层 + L1/L2/L3 三层沙箱 \\
fuxi-memory & 2\,643 & 长期记忆：甲骨 / 河图洛书 / 身份卡 \\
fuxi-agent-codex & 2\,197 & Codex CLI 门客适配器（spawn-per-dispatch） \\
fuxi-scheduler & 2\,146 & cron / once / fs / webhook 触发器调度 \\
fuxi-events & 2\,135 & EventBus + SQLite WAL + 历史回放 \\
fuxi-firehose & 2\,135 & WebSocket + SSE + REST + TUI 四端输出 \\
fuxi-core & 2\,102 & 核心 trait：Agent / Event / Task / Workspace \\
fuxi-a2a & 1\,321 & A2A v1.0 wire 协议 + JSON-RPC server + tests \\
fuxi-skills & 1\,116 & 角色玉牒加载器 + 招贤流程 \\
\midrule
\textbf{总计} & \textbf{82\,515} & \\
\bottomrule
\end{tabular}
\end{table}

依赖图的设计原则是"单向依赖、无循环"。最底层是 \texttt{fuxi-core}，只声明 trait 与基础类型（\texttt{Agent}、\texttt{Workspace}、\texttt{Event}、\texttt{Task}），不依赖任何其他 fuxi 内部 crate；中间层是 \texttt{fuxi-events}、\texttt{fuxi-a2a}、\texttt{fuxi-workspace} 等基础设施，单向依赖 \texttt{fuxi-core}；编排层 \texttt{fuxi-orchestrator} 依赖中间层，并定义 \texttt{RecallSink}、\texttt{DistEnqueuer} 等扩展点 trait；适配器层 \texttt{fuxi-agent-cc}、\texttt{fuxi-agent-codex} 实现 \texttt{fuxi-core::Agent}；最顶层是 \texttt{fuxi-cli}，把所有 crate 装配起来并实现各扩展点 trait。值得指出的是，\texttt{fuxi-memory} 需要在某些路径上调用编排层逻辑（例如在 deliverable 入库时触发 oracle 抽取），但它自身不依赖 \texttt{fuxi-orchestrator}——而是反向定义抽取器 trait 由 \texttt{fuxi-cli} 实现并注入。这种"trait 在低层、impl 在高层"的反向依赖模式是规避循环依赖的标准 Rust 工程惯用法，伏羲在 \texttt{RecallSink}、\texttt{DistEnqueuer}、\texttt{NodeLoadProvider} 等扩展点上反复使用了它。

值得展开说明的是 \texttt{fuxi-cli} 占了整个 workspace 近 44\% 的代码量。这反映了一个工程现实：协议、抽象与基础设施虽然概念上是"重头戏"，但把它们装配成一个能跑起来、能被用户实测的产品需要大量的胶水代码。\texttt{fuxi-cli} 内部进一步分成 REPL 子模块（用户在终端的交互入口）、daemon 子模块（后台常驻服务）、dist 子模块（v2 跨节点 controller）、im 子模块（PWA 后端集成）、xuannv\_cmd 子模块（玄女生命周期管理）、ipc 子模块（REPL 与 daemon 之间的 Unix socket 通讯）等十余个垂直功能。这种"集成层占大头"的代码分布在中大型 Rust 项目里很常见，它也解释了为什么单 binary crate 内的 lint 与测试覆盖反而需要更高的纪律——基础设施层的 bug 通常在编译期就能暴露，但集成层的 bug 往往要靠端到端测试才能复现。

13 个 crate 的命名遵循统一的 \texttt{fuxi-*} 前缀。早期项目曾用 \texttt{xihe-*} 前缀（"羲和"的拼写变体），在 2026-04-19 项目正名为伏羲后所有 crate 名都迁到了新前缀，旧路径通过文件系统软链接保持向后兼容（\texttt{/Users/e0\_7/xihe} 仍指向 \texttt{/Users/e0\_7/fuxi}）。这种"文件系统层兼容、代码层硬迁移"的做法在重命名重大模块时是个常见的折中——既不破坏已有终端会话与脚本，又强制新代码立即采用新约定，避免 \texttt{xihe-} 与 \texttt{fuxi-} 两套命名长期并存的认知负担。

\subsection{事件数据结构与 Lamport 因果序}

事件是伏羲所有跨模块、跨进程通讯的统一载体。一个事件由两部分组成：\texttt{EventMeta} 携带通用元数据，\texttt{EventKind} 是带 50+ 变体的标签联合（tagged union），描述事件具体类型。代码~\ref{lst:event-types} 节选自 \texttt{crates/fuxi-core/src/event.rs}，给出 \texttt{Event} 与 \texttt{EventKind} 的核心定义。

\begin{lstlisting}[language=Rust, caption={Event 与 EventKind 核心定义（节选自 fuxi-core/src/event.rs）}, label={lst:event-types}]
pub struct Event {
    pub meta: EventMeta,   // id / at / session / agent / task / source_node_id
    pub kind: EventKind,
}

#[derive(Debug, Clone, Serialize, Deserialize)]
#[serde(tag = "type", rename_all = "snake_case")]
pub enum EventKind {
    AgentSpawning { role: String, cli: String },
    AgentReady { endpoint: String },
    AgentDead { cause: String },
    TaskCreated { title: String, description: String },
    TaskDispatched { target: Option<AgentId> },
    TaskStateChanged { from: TaskState, to: TaskState },
    AgentResponded { text: String },
    ToolCallStarted { tool: String, args: serde_json::Value },
    ToolCallFinished { tool: String, result: serde_json::Value },
    DeliverableProduced { deliverable_kind: DeliverableKind,
                           files: Vec<DeliverableFileMeta> },
    WorkspaceCreated { workspace_id: WorkspaceId, layer: WorkspaceLayer },
    WorkerHeartbeat { node_id: String, inflight: u32 },
    TriggerFired { spec: TriggerSpec, cause: FireCause },
    Custom { label: String, payload: serde_json::Value },
    // ... 共 50+ 变体
}
\end{lstlisting}

\texttt{\#[serde(tag = "type", rename\_all = "snake\_case")]} 这条 \texttt{serde} 属性是设计的关键：序列化时每个变体会被翻译成一个 JSON 对象，其中 \texttt{type} 字段为 \texttt{"agent\_spawning"} / \texttt{"task\_created"} 等小写字符串，其余字段直接拍平到对象顶层。相比无 \texttt{tag} 的 \texttt{untagged} 联合，这种显式标签方式有两个不可替代的优点：一是反序列化无歧义，不会因为两个变体字段重合而出现回退优先级错乱；二是新增变体时所有 \texttt{match ev.kind} 处都会编译失败，强制开发者更新所有消费路径——这正是 Rust enum 与 \texttt{serde} 联合发挥类型安全价值的典型场景。在伏羲过去一年的演进中，\texttt{EventKind} 从最初的 8 个变体扩展到 50+ 变体，每次新增都依赖编译器穷举检查找出所有需要更新的 firehose 渲染、SQL 持久化与 TUI 摘要逻辑，这种"编译期保证不落下任何消费者"的特性是字符串 tag 方案无法企及的。

事件的持久化由 \texttt{EventStore} 负责。底层是单文件 SQLite 库，开启 WAL 模式以支持并发读写。表结构定义在 \texttt{crates/fuxi-events/migrations/0001\_init.sql}，并通过 \texttt{include\_str!} 在编译期内嵌到 \texttt{store.rs}，这样即使是 \texttt{:memory:} 测试库也能用同一份 schema 初始化，避免运行时路径依赖。代码~\ref{lst:event-schema} 给出表结构。

\begin{lstlisting}[language=SQL, caption={事件审计日志表结构（fuxi-events/migrations/0001\_init.sql）}, label={lst:event-schema}]
CREATE TABLE IF NOT EXISTS events (
    id            TEXT PRIMARY KEY,    -- uuid
    at            TEXT NOT NULL,       -- ISO 8601，时间游标
    session       TEXT,
    agent         TEXT,
    task          TEXT,
    kind_tag      TEXT NOT NULL,       -- serde "type" tag，过滤索引用
    payload       TEXT NOT NULL,       -- 完整 Event JSON
    -- 同秒内插入需要 tiebreaker；rowid 自增即可还原插入顺序
    recorded_at   TEXT NOT NULL DEFAULT (strftime('%Y-%m-%dT%H:%M:%f', 'now'))
);
CREATE INDEX IF NOT EXISTS idx_events_at ON events(at);
CREATE INDEX IF NOT EXISTS idx_events_agent ON events(agent);
CREATE INDEX IF NOT EXISTS idx_events_task ON events(task);
\end{lstlisting}

\texttt{kind\_tag} 字段独立索引，使按事件类型过滤（"找出所有 \texttt{AgentSpawning}"）退化为 O(log n) 的索引扫描；\texttt{recorded\_at} 用毫秒精度的字符串配合自增 rowid 作 tiebreaker，保证同秒插入的事件仍有稳定的全序。\texttt{at} 字段使用 ISO 8601 字符串而非 unix timestamp 整数，是因为字符串可以直接在 SQL 里按字典序排序得到时间序，且人类可读——\texttt{sqlite3 events.db "SELECT at, kind\_tag FROM events WHERE task = ?"} 这种 ad-hoc 排查查询在事故复盘时极为高频，整数 timestamp 还得人脑或脚本反算成时间反而拖慢节奏。\texttt{payload} 字段存放整条 \texttt{Event} 序列化为 JSON 的全文，看似冗余（部分字段已抽出到独立列），但保证了 schema 演进时的零破坏迁移：未来如果加新字段或新枚举变体，只要 serde 反序列化兼容旧 \texttt{payload}，老库无需 \texttt{ALTER TABLE} 即可继续读取。这种"独立列做索引、整体 JSON 做真相"的混合存储是事件溯源系统的常见模式，类似 Kafka 把消息体当 opaque blob、把分区/offset/key 抽到 metadata 层做索引\cite{kreps2011kafka}。

事件之间的因果关系采用 Lamport 逻辑时钟\cite{lamport1978time}。设事件 $a$ 是本地新事件、$m$ 是从外部节点收到的事件、$C(\cdot)$ 表示某事件携带的逻辑时钟，则更新规则为：
\begin{equation}
\begin{aligned}
C(a) &= C_{\text{local}} + 1, \\
C(\text{recv}(m)) &= \max(C_{\text{local}}, C(m)) + 1.
\end{aligned}
\label{eq:lamport}
\end{equation}
式~\eqref{eq:lamport} 蕴含的"如果 $e_1$ 因果先于 $e_2$，则 $C(e_1) < C(e_2)$"性质，使得伏羲在跨节点回放、IM 端历史拉取、TUI 滚动等场景下都能给出一致的事件全序。\texttt{EventMeta.source\_node\_id} 字段用于在分布式部署中区分事件起源节点：当 dist controller 把 worker 端事件 republish 到 home 节点的 EventBus 时，会保留原 \texttt{source\_node\_id}，使下游消费者可以判断"这是远端起源"还是"本地原生"。

需要说明的是，伏羲在多数应用场景下并不强依赖 Lamport 时钟的全序性质——单进程内事件的真实物理时间戳已经足够精确（同进程内 \texttt{Instant::now} 单调递增），多进程跨节点的事件即便偶有顺序漂移也能被业务语义弥补（例如 \texttt{TaskCompleted} 必然在 \texttt{TaskCreated} 之后）。Lamport 时钟在这里的真正价值是给"分布式断言"提供基础设施：当用户问"为什么 worker A 上的 task 完成事件比 worker B 上的开始事件还早"时，逻辑时钟可以明确指出"它们之间没有因果关系"——这一回答比"机器时钟漂移"要可靠得多，特别是在 NTP 同步偶尔失灵、个人开发者机器经常休眠唤醒的实际部署里。

\texttt{EventBus} 的发布路径是整个事件系统的性能关键点。一次 \texttt{publish} 必须满足两个看似矛盾的约束：（1）发布者绝不被阻塞——任何阻塞都可能让 dispatch pump 回压，进而堆积上游门客的 stdout 缓冲区；（2）事件不允许丢失——审计日志一旦缺一条就破坏了"SQLite 是单一真相源"这条公理。代码~\ref{lst:eventbus-publish} 给出实际实现。

\begin{lstlisting}[language=Rust, caption={EventBus::publish 非阻塞实现（fuxi-events/src/bus.rs）}, label={lst:eventbus-publish}]
pub fn publish(&self, ev: Event) -> Result<()> {
    // 1) 广播扇出——0 receiver 不算错。
    let _ = self.inner.broadcast_tx.send(ev.clone());

    // 2) 写入：try_send 非阻塞塞进 mpsc。
    match self.inner.writer_tx.try_send(ev) {
        Ok(_) => {
            self.inner.writer_pending.fetch_add(1, Ordering::Relaxed);
        }
        Err(mpsc::error::TrySendError::Full(dropped)) => {
            // 公理：原始 publish 不能被丢——宁可阻塞背景任务也要塞进去。
            let pending = self.inner.writer_pending.load(Ordering::Relaxed);
            if pending >= self.inner.cfg.lag_threshold {
                self.maybe_emit_lag_sentinel(pending);   // 哨兵告警
            }
            let writer_tx = self.inner.writer_tx.clone();
            let pending_counter = self.inner.writer_pending.clone();
            tokio::spawn(async move {
                if writer_tx.send(dropped).await.is_ok() {
                    pending_counter.fetch_add(1, Ordering::Relaxed);
                }
            });
        }
        Err(mpsc::error::TrySendError::Closed(_)) => {
            return Err(Error::WriterClosed);
        }
    }
    Ok(())
}
\end{lstlisting}

设计要点有三：第一，broadcast 通道与 mpsc 持久化通道完全分离——broadcast 负责实时扇出（订阅者各取所需，零拷贝），mpsc 负责串行落库（FIFO 入 SQLite）。这两条路径互不阻塞，慢订阅者拖累不到快订阅者，写库拥塞也阻塞不了实时推送。第二，\texttt{try\_send} 满了之后不丢消息，而是把当前事件 \texttt{spawn} 到一个后台 task 让它去 \texttt{send().await}——这个后台 task 可以阻塞，但不会拖到原 \texttt{publish} 调用方。第三，当 \texttt{writer\_pending} 计数超过阈值时，\texttt{maybe\_emit\_lag\_sentinel} 会主动发一条 \texttt{Custom \{ label: "event\_store\_lagged", payload \}} 事件到 broadcast——这是一种"自我观测"的可观测性手段，订阅者（如 firehose TUI）能直接看到 backpressure 信号而不需要外部监控。

订阅端必须容忍 broadcast 的 \texttt{Lagged} 错误。tokio 的 \texttt{broadcast::Receiver} 在订阅者跟不上节奏时会返回 \texttt{RecvError::Lagged(n)}，含义是"你已经落掉了 N 条事件，下次接收将从最新可见位置继续"。伏羲的策略是：把 \texttt{Lagged} 当 warn 级日志记一下，继续接收新事件，不让单次 lag 杀死整个订阅链。完整历史另走 \texttt{replay(cursor, live\_tail=true)} 从 SQLite 拉，订阅流只承担"实时推送"职责。这种双路径分工把"实时性"与"完整性"解耦到了不同代码路径，各自取最优实现。

写库这一侧还有一个工程细节需要处理：SQLite WAL 模式在高并发写入时偶发 \texttt{SQLITE\_BUSY} 错误。\texttt{EventStore::append} 通过 sqlx 的 \texttt{busy\_timeout(5s)} 已能消化 99\% 的瞬时冲突，但在压力测试里仍会有零星 BUSY 漏出来。代码~\ref{lst:eventstore-append} 给出三次指数退避的兜底实现。

\begin{lstlisting}[language=Rust, caption={EventStore::append 三次指数退避（fuxi-events/src/store.rs）}, label={lst:eventstore-append}]
pub async fn append(&self, ev: &Event) -> Result<()> {
    let id = ev.meta.id.to_string();
    // ... 拆 meta / 序列化 payload 略
    let mut attempt: u32 = 0;
    loop {
        let res = sqlx::query(
            "INSERT INTO events (id, at, session, agent, task,
                                  kind_tag, payload) \
             VALUES (?1, ?2, ?3, ?4, ?5, ?6, ?7)",
        ).bind(&id).bind(&at).bind(&session).bind(&agent)
         .bind(&task).bind(kind_tag).bind(&payload)
         .execute(&self.pool).await;
        match res {
            Ok(_) => return Ok(()),
            Err(sqlx::Error::Database(db))
                if is_busy(db.as_ref()) && attempt < 3 =>
            {
                attempt = attempt.saturating_add(1);
                tokio::time::sleep(Duration::from_millis(
                    50_u64.saturating_mul(attempt as u64),
                )).await;     // 50ms / 100ms / 150ms 退避
            }
            Err(e) => return Err(e.into()),
        }
    }
}
\end{lstlisting}

这段代码是"层层防御"思路的典型。第一层 \texttt{busy\_timeout} 在 SQLite 内部等锁，吞掉绝大多数瞬时冲突；第二层应用层的 3 次指数退避兜底（50ms / 100ms / 150ms），把仍然漏出的偶发错误压到统计零；第三层是 \texttt{publish} 中的后台 spawn 转交，确保即便 append 真的失败几次，原 \texttt{publish} 调用方仍不被阻塞。三层防御加在一起，伏羲在过去半年实测中没有出现过事件审计日志缺失的事故。值得指出的是，\texttt{:memory:} 测试库有一个独特陷阱——同一 \texttt{:memory:} URI 的不同 connection 在 SQLite 内是不同的库，多连接写入后读不到。\texttt{EventStore::new} 检测到 \texttt{:memory:} 时强制 \texttt{max\_connections = 1} 来规避这个陷阱，避免单元测试里出现"明明 append 了却 replay 不出来"的诡异行为。

\subsection{任务派发流程实现}

任务派发是编排层（\texttt{fuxi-orchestrator}）最核心的代码路径。当用户在 IM 端或 TUI 输入一句话后，玄女判断是否需要派单，若需要则调用 \texttt{Fuxi::dispatch} 或 \texttt{Fuxi::dispatch\_to\_any}。派发的核心责任有四项：（1）按 \texttt{pinned\_node} / \texttt{required\_tags} / \texttt{project.host\_nodes} 决策本地或分布式路径；（2）对本地路径，调 \texttt{Agent::dispatch} 拿到 \texttt{mpsc::Receiver<Event>}；（3）启动 dispatch pump 把私有 receiver 中的事件 republish 到全局 EventBus；（4）在 pump 终止时清理状态、发 \texttt{TaskCompleted}/\texttt{TaskCancelled} 事件。整体流程见后文 \autoref{lst:dispatch-main} 与代码内注释。

代码~\ref{lst:claim-idle} 是 \texttt{Shelf} 的 \texttt{claim\_idle\_by\_role}——一个原子的"找空闲门客并占用"操作，是规避 TOCTOU 双派 bug 的关键。

\begin{lstlisting}[language=Rust, caption={原子地按角色 claim 一个空闲门客（fuxi-orchestrator/src/registry.rs）}, label={lst:claim-idle}]
pub async fn claim_idle_by_role(&self, role: &str) -> Option<AgentId> {
    let mut guard = self.inner.write().await;
    let pick = guard
        .iter_mut()
        .find(|(_, e)| e.status == ShelfStatus::Idle
              && e.card.profile.role == role);
    if let Some((id, e)) = pick {
        e.status = ShelfStatus::Busy;
        // claim 等价于"占用"：必须清 idle_since，否则 GC 仍把这只门客
        // 算作 idle 超时，竞态会把已 busy 的门客误杀。
        e.idle_since = None;
        Some(*id)
    } else {
        None
    }
}
\end{lstlisting}

这段代码看起来朴素但藏着两个并发陷阱。第一个是 TOCTOU（Time-of-Check-Time-of-Use）：早期版本把 \texttt{find\_idle\_by\_role} 与后续 \texttt{set\_status(Busy)} 拆成两步，结果两个并发 dispatch 都 \texttt{find} 到同一只门客，再分别 \texttt{set\_status}，导致同一个门客被双派而事件流相互覆盖；现在 \texttt{find + status update} 全部在 \texttt{write().await} 锁内完成，自然原子。第二个是 idle GC 的清扫线程：另一条独立线程会扫 \texttt{idle\_since} 字段超过阈值的门客并 \texttt{shutdown\_agent}，如果 claim 时不清掉 \texttt{idle\_since}，已 busy 的门客仍被 GC 视作"长时间空闲"而误杀；显式置 \texttt{None} 才能让 GC 与 claim 两条路径不冲突。这种"原子性"与"路径无副作用"的细节，是新接手系统的开发者最容易踩进去的坑——它们看上去都是"小动作"，实际上每一处都对应过一次线上事故。整个 \texttt{registry.rs} 的 1\,843 行代码里至少有 12 处类似的"显式 \texttt{None} / 显式锁内更新 / 显式 invariant 检查"注释，每一条都标记了被修复的 bug 编号。

代码~\ref{lst:dispatch-main} 是 \texttt{Fuxi::dispatch} 的本地路径主循环（删去了 v2 分布式分支与 sentinel addendum 拼装的细节，保留控制流主干）。

\begin{lstlisting}[language=Rust, caption={dispatch 本地路径主循环（节选自 fuxi-orchestrator/src/fuxi.rs）}, label={lst:dispatch-main}]
pub async fn dispatch(&self, agent_id: AgentId, task: Task) -> Result<()> {
    // 1. v2 自动 pin：未显式 pin 但 task 关联 project 时，按
    //    project.host_nodes 选最闲节点 auto-pin。
    let mut task = task;
    if let Some(picked) = self.auto_pin_from_project(&task).await {
        task.pinned_node = Some(picked);
    }

    // 2. 路由决策：pinned_node 或 required_tags 非空 -> 走 dist enqueue。
    let needs_dist = task.pinned_node.is_some()
                     || !task.required_tags.is_empty();
    if needs_dist {
        if let Some(enq) = self.dist_enqueuer.read().await.clone() {
            // 远端 worker 跑：补发 TaskCreated 让 home aggregate 拿 title
            self.bus.publish(Event {
                meta: EventMeta::for_task(task.id),
                kind: EventKind::TaskCreated {
                    title: task.title.clone(),
                    description: task.description.clone(),
                },
            })?;
            enq.enqueue(&task, DistEnqueueOptions {
                system_prompt: self.compose_role_prompt(agent_id).await,
                task_id: Some(task.id.to_string()),
                role: Some(self.shelf.role_of(agent_id).await?),
            }).await?;
            return Ok(());
        }
        warn!(task = %task.id, "needs_dist but enqueuer not injected; fallback");
    }

    // 3. 本地路径：取 agent，启动 receiver，spawn pump task。
    let agent = self.shelf.get_agent(agent_id).await
        .ok_or(Error::AgentNotFound)?;
    let task_id = task.id;
    let rx = agent.dispatch(task.clone()).await?;
    let bus = self.bus.clone();

    tokio::spawn(async move {
        let mut stream = ReceiverStream::new(rx);
        let mut saw_terminal = false;
        loop {
            // 终态后给 50ms grace window，让晚到的 deliverable 事件
            // 也能 republish; 超时再退。这是 M2.1 修过的尾包丢失 bug。
            let next = tokio::time::timeout(
                Duration::from_millis(50),
                stream.next(),
            ).await;
            match next {
                Ok(Some(mut ev)) => {
                    ev.meta.task = Some(task_id);
                    let _ = bus.publish(ev.clone());
                    if ev.kind.is_terminal() {
                        saw_terminal = true;
                    }
                }
                Ok(None) => break,             // channel closed
                Err(_) if saw_terminal => break, // grace 超时
                Err(_) => continue,            // 还没看到终态，继续等
            }
        }
    });
    Ok(())
}
\end{lstlisting}

本地路径的核心控制流可以归纳为五步：取 agent 句柄、调 \texttt{Agent::dispatch} 拿 \texttt{mpsc::Receiver<Event>}、spawn 后台 pump task、pump task 内补 \texttt{task\_id} 并 republish、终态 + grace window 后退出。这五步看似简单，但每一步都对应一个具体的 invariant：取句柄前先确认 agent 仍在 shelf 中、dispatch 必须返回 receiver（不允许返 stream 或 future）、pump task 必须能在 shelf 关停时优雅退出、补 \texttt{task\_id} 是因为 agent adapter 不知道 platform 层的 task 概念、grace window 是为了不丢尾包。把这些 invariant 写在代码注释里而非孤立的设计文档里，未来重构时新开发者一眼就能看到"为什么不能省这一步"。

这段代码体现了 publish-then-pump 模式：编排层对外只暴露 \texttt{dispatch(agent\_id, task) -> Result<()>}，调用方不关心事件如何流转；门客返回的 \texttt{mpsc::Receiver<Event>} 是私有通道，只在 pump task 内部消费，每条事件被补全 \texttt{task\_id} 后立即 \texttt{bus.publish()} 到全局通道。这样一来，订阅 EventBus 的所有下游（firehose TUI、IM WebSocket、SCM 提取器、玄女自身）都能看到完整事件流，而不需要分别去敲门客。\texttt{is\_terminal()} 后保留 50ms grace window 是 M2.1 实测发现的 bug 修复——cc 在发出 \texttt{ResultSuccess} 后还会异步 flush 几条 \texttt{DeliverableProduced} 事件，立即 break 会让它们被丢，给 50ms 缓冲后再退则全部能 republish。

\texttt{auto\_pin\_from\_project} 这个步骤在 v2 跨节点版本里加入。它的语义是："如果 task 关联到某个 \texttt{Project} 且未显式 pinned\_node，就按 \texttt{project.host\_nodes} 列表选最闲的节点 auto-pin"。实现上调 \texttt{NodeLoadProvider}（CLI 装配时注入的 \texttt{Arc<dyn NodeLoadProvider>}）取每个候选节点的 \texttt{saturation = inflight / max\_concurrency}，pick 最低 saturation 的；候选全离线时返 None 退化到本地 spawn 路径，保证向后兼容。这条路径有一个反复被绕进去的边界条件：用户显式 \texttt{@<node>} pin 时绝不能被 auto\_pin 改写——\texttt{auto\_pin\_from\_project} 入口先看 \texttt{task.pinned\_node.is\_some()}，已 pin 直接返 None；调用方再加 \texttt{is\_none()} 守卫双保险。这种「双重防御」是多次实测后沉淀的规约：单层守卫一旦在重构中被误删，就会出现「用户显式 pin 在 mac-mini 上的任务最终跑到 home 节点」的回归现象。

dispatch 还有一个不太显眼但生产关键的细节：\texttt{Fuxi::shutdown\_agent} 必须豁免玄女本人。idle GC 任务会扫描超过 10 分钟未活动的门客并 \texttt{shutdown\_agent} 释放资源，但玄女对 GC 而言是普通 agent——如果不豁免，长时间没人输入时 GC 会把玄女自己关掉，整个 TUI 立刻无响应。\texttt{shutdown\_agent} 入口比对 \texttt{xuannv\_id()}，命中就 warn 后返 \texttt{Ok(())} 静默 noop。任何新增的 shutdown 路径（\texttt{fuxi kill --id} / 将来的 worker rebalance）都从这条豁免方法走，不绕旁路；唯一例外是 task \#8 上下文交接专用的 \texttt{shutdown\_xuannv\_for\_handoff}，命名上故意带 \texttt{\_for\_handoff} 让 grep 一眼看清"不要在别处调用"。此类细粒度的约束在代码层强制成立，是系统稳定性的来源之一。

玄女如何感知 worker 端的状态变化？答案是 \texttt{SystemEventBridge}。它是一个独立的 tokio task，订阅 EventBus，对 \texttt{AgentDead}/\texttt{TriggerFired}/\texttt{AgentRequestReview}/\texttt{TaskStateChanged} 这些"系统事件"做翻译，转成"系统提示文本"通过 \texttt{Intervener::intervene\_system\_origin} 喂给玄女。代码~\ref{lst:bridge-spawn} 给出 spawn 入口。

\begin{lstlisting}[language=Rust, caption={SystemEventBridge spawn 入口（fuxi-orchestrator/src/bridge.rs）}, label={lst:bridge-spawn}]
fn spawn_inner(
    intervener: Arc<dyn Intervener>,
    bus: EventBus,
    xuannv_id: AgentId,
    trigger_lookup: Arc<dyn TriggerLookup>,
    backoff_ms: &'static [u64],
) -> JoinHandle<()> {
    let mut sub = bus.subscribe();
    let bus_for_handler = bus.clone();
    tokio::spawn(async move {
        while let Some(item) = sub.next().await {
            let Ok(ev) = item else { continue; };  // Lagged 已被 filter
            handle_event(&*intervener, &*trigger_lookup, xuannv_id,
                         &bus_for_handler, backoff_ms, ev).await;
        }
        debug!("SystemEventBridge: 订阅流结束，退出");
    })
}
\end{lstlisting}

这条桥是公理 \#3"真实时不轮询"的具体落地——玄女完全不查询任何状态，所有"知情权"都通过 EventBus 推送实现。把"订阅"与"翻译"职责放在独立 crate 内，使得"如果将来想换玄女实现（比如换成另一个本地大模型）"只要重新实现 \texttt{Intervener} trait，这条桥代码可以原样复用。

bridge 内部还有一处关键工程细节：\texttt{handle\_event} 把 \texttt{EventKind} 翻译成「系统提示文本」时会做去重与节流。例如 worker 节点 inflight 较高时会高频发出 \texttt{WorkerHeartbeat}，若每条都翻译成 system prompt 喂给玄女，会让她被心跳事件刷屏而无法处理其他工作；bridge 内部维护一份「上次喂给玄女的同类事件时间戳」映射，只有跨过抑制窗口的事件才真正被翻译。\texttt{ReviewRequestTimeout} 这种关键事件还有 backoff 重试——首次翻译失败时按 \texttt{REVIEW\_RETRY\_BACKOFF\_MS}（默认 \texttt{[100, 300, 1000]} 三次）重试，避免玄女正在处理别的事时一次失败就丢掉重要通知。这种"事件路由 + 节流 + 重试"的组合在 fuxi 早期版本里散落在多处代码中，M2 重构后集中到 bridge crate 内单元测试覆盖率达 80\%，是降低系统运行复杂度的关键一步。

\subsection{A2A 协议从零实现（论文核心贡献）}

A2A（Agent-to-Agent）是 Google 在 2025 年提出的开放 agent 互通协议\cite{a2aproject2025a2a}，其规范级别等价于 MCP 在工具调用领域的位置\cite{anthropic2024mcp}。这两个协议构成了 LLM 时代 agent 生态的两条主轴：MCP 解决"agent 如何调用工具"，A2A 解决"agent 如何调用其他 agent"。然而截至本文写作时，A2A 在 Rust 生态没有任何官方或第三方实现——Google 仅发布了 Python 与 Go 参考实现，社区里另有 Java/TypeScript 民间移植，而其他企业级实现往往与具体框架（如 SWE-Agent\cite{yang2024sweagent}、OpenHands\cite{wang2024openhands}）深度耦合，不能直接拿来即用。伏羲为了让玄女能与未来的异构 agent（其他公司发布的、运行在其他主机上的 cc/gemini/codex 实例）互联，必须从零实现一个完整的 A2A 1.0 客户端与服务端。这一节是本文工程贡献的核心部分。

需要先澄清一个常见误解：A2A 协议并不规定"agent 如何思考"或"agent 用什么模型"，它只规定 agent 之间如何握手、如何传递任务、如何回报结果。这种"协议关心通讯、不关心实现"的设计哲学跟 HTTP 类似——HTTP 规范不规定网页用什么框架渲染，只规定字节流的形态。A2A 这种设计让伏羲可以把 cc / codex / 未来的 gemini-cli 通通包装成 A2A endpoint 后由玄女统一调度，不必为每个 agent 实现单独的协议适配器。

设计这个 crate 时面临的第一个根本问题是"实现到什么粒度"。A2A 1.0 规范涵盖了 agent 自描述、任务生命周期、流式推送、artifact 分块上传、push notification、双向 TLS、版本协商等若干主题，全部实现工程量巨大。我们采取了"覆盖 fuxi 场景所需闭环、明确简化项"的务实策略：覆盖 5 个核心 RPC 方法（\texttt{agent/getCard} / \texttt{tasks/send} / \texttt{tasks/sendSubscribe} / \texttt{tasks/get} / \texttt{tasks/cancel}）及对应的 \texttt{TaskState} 状态机（含 fuxi 自创的 \texttt{InputRequired}）、\texttt{Message} / \texttt{Part} / \texttt{Artifact} 层次结构以及 SSE 流式分块；明确简化掉 push notification（fuxi 内网部署，无 webhook 基础设施）、双向 TLS（信任边界在进程而非网络）、自动重连（上层 \texttt{A2ARouter} 负责）、版本协商（假设对端实现兼容 A2A 1.0）这四类位于 A2A v1.0 规范周边的能力。这一覆盖范围对应 A2A v1.0 的核心 RPC 闭环——足以与任意 A2A 1.0 兼容方完成 agent discovery、task send、task stream、task query、task cancel 五条主路径的互通；crate 体积也因此控制在 1\,063 行核心代码，便于审阅。

\subsubsection{Crate 模块拆分与协议规范差异}

\texttt{crates/fuxi-a2a} 拆分为 7 个文件：\texttt{wire}（458 行，规范数据结构）、\texttt{jsonrpc}（77 行，envelope）、\texttt{sse}（45 行，事件流抽象）、\texttt{server}（196 行，Axum HTTP 端点）、\texttt{client}（203 行，Reqwest 客户端）、\texttt{error}（52 行，错误类型）、\texttt{lib.rs}（32 行，re-export），加上端到端测试 \texttt{tests/roundtrip.rs} 258 行，总计 1\,321 行。这个比例（核心库 1063 行、测试 258 行、文档注释占 25\%）反映了一个"协议库应有的密度"——逻辑紧凑但每个公开 API 都有完整文档。

实现上与 A2A 1.0 规范保持互通的同时做了以下五处差异化处理：

\begin{table}[H]
\centering
\caption{fuxi-a2a 与 A2A 1.0 规范的字段差异}
\label{tab:a2a-diff}
\begin{tabular}{lll}
\toprule
\textbf{项} & \textbf{A2A 1.0} & \textbf{fuxi 实现} \\
\midrule
字段命名 wire & camelCase & camelCase（serde rename） \\
\texttt{url} 字段 & 可选 & 必需（fuxi 需立即知 endpoint） \\
\texttt{streaming} & 可选 & 总是 true（统一走 SSE） \\
\texttt{pushNotifications} & 可选 & 总是 false（v1 暂无 webhook） \\
\texttt{InputRequired} & 无 & 新增（fuxi 创新，人工介入点） \\
\bottomrule
\end{tabular}
\end{table}

\texttt{InputRequired} 是最重要的一项扩展。Google 规范中 \texttt{TaskState} 枚举只有 \texttt{Submitted/Working/Completed/Failed/Canceled} 五态，全部是"agent 自治"语义。但伏羲的工作模式是"人在回路"——门客在某些决策点必须等用户确认（例如 reviewer 看完代码后是否合并），因此显式增加 \texttt{InputRequired = "input-required"} 状态，让对端编排器（玄女）能感知"这个 task 不是失败，而是等人工"。代码~\ref{lst:taskstate} 给出最终定义。

\begin{lstlisting}[language=Rust, caption={TaskState wire 形态（fuxi-a2a/src/wire.rs）}, label={lst:taskstate}]
/// 状态机 on wire——kebab-case。
///
/// `InputRequired` 映射到 `"input-required"`：把人工介入显式地暴露给
/// 对端编排器（玄女），让她能决策是否转交主对话权。
#[derive(Debug, Clone, Copy, PartialEq, Eq, Serialize, Deserialize)]
#[serde(rename_all = "kebab-case")]
pub enum TaskState {
    Submitted,
    Working,
    InputRequired,
    Completed,
    Failed,
    Canceled,
}
\end{lstlisting}

\texttt{AgentCard} 是 A2A 协议中"agent 自描述"的载体，每个对外暴露 A2A 端点的 agent 必须能通过 \texttt{agent/getCard} 方法返回一份 \texttt{AgentCard} JSON。代码~\ref{lst:agentcard} 给出 fuxi 的 wire 形态。

\begin{lstlisting}[language=Rust, caption={A2A AgentCard wire 定义（fuxi-a2a/src/wire.rs）}, label={lst:agentcard}]
#[derive(Debug, Clone, PartialEq, Eq, Serialize, Deserialize)]
pub struct AgentCard {
    pub name: String,                    // "luban-1"
    pub description: String,             // 来自 system_prompt 首行（截 200 字）
    pub version: String,                 // env!("CARGO_PKG_VERSION")
    pub url: String,                     // 本 agent 的 A2A POST 端点
    pub capabilities: AgentCapabilities, // streaming / pushNotifications
    pub skills: Vec<AgentSkill>,         // 从 fuxi tags 平铺
}

#[derive(Debug, Clone, Copy, PartialEq, Eq, Serialize, Deserialize)]
pub struct AgentCapabilities {
    pub streaming: bool,
    #[serde(rename = "pushNotifications")]
    pub push_notifications: bool,
}
\end{lstlisting}

注意 \texttt{pushNotifications} 字段——A2A 规范在 wire 上用 camelCase，而 idiomatic Rust 是 snake\_case，因此用 \texttt{\#[serde(rename = "pushNotifications")]} 在序列化层做翻译。这种"一字段两形态"在整个 wire.rs 里反复出现，例如 \texttt{inputModes}/\texttt{outputModes}/\texttt{historyLength}/\texttt{acceptedOutputModes} 等。Rust 侧代码全用 snake\_case，wire 上自动按规范输出 camelCase，开发者既享受 idiomatic 体验又不破坏跨语言互通。这种命名翻译看似琐碎，但在大规模协议实现里是大量代码行的来源——A2A wire types 共有 11 个字段需要 \texttt{rename}、6 个枚举需要 \texttt{rename\_all = "kebab-case"}，全部由 serde 属性宏在编译期处理，运行时零开销。

内部 \texttt{fuxi\_core::AgentCard}（带 \texttt{AgentId}、\texttt{AgentStatus} 等运行时字段）到对外 \texttt{wire::AgentCard} 的转换是有损的——\texttt{id} 与 \texttt{status} 都不应暴露给外部。fuxi 用 \texttt{From} trait 在边界上原子转换，禁止其他模块手抄字段（M3.4 约束）。\texttt{description} 取 \texttt{system\_prompt} 第一行截 200 字符；\texttt{capabilities.streaming = true}（fuxi 统一 SSE）；\texttt{skills} 从 \texttt{profile.tags} 平铺成 stub，每个 tag 一个 \texttt{AgentSkill}。这条 \texttt{From} 实现的关键是"信息单向流动"——只允许从 core 到 wire，wire 类型本身没有也不应有 \texttt{Into<core::AgentCard>}，因为对外 wire 缺少 \texttt{AgentId} 与 \texttt{AgentStatus} 这些内部 invariant 必需的字段。这种"两个相似但不相等的类型"的边界设计在跨进程协议中很常见，强行用一个共享 struct 兼容内外两种语义往往会带来"内部字段意外泄漏到 wire""wire 上的 Optional 字段在内部被错误依赖"等微妙 bug。\texttt{Part} 类型的设计同样体现了这种严谨——伏羲选择 \texttt{\#[serde(tag = "type", rename\_all = "kebab-case")]} 而非 \texttt{untagged}，确保 wire 上每个 \texttt{Part} 都显式标记类型（\texttt{"type": "text"} / \texttt{"type": "data"} / \texttt{"type": "file"}），避免"文本恰好长得像 JSON"或"data 字段碰巧含 \texttt{name} 键"等场景下的反序列化优先级歧义。这一处微小的协议细节，在生产系统里能避免大量"为什么这条 message 被识别为 file part"式的诡异 bug。

\subsubsection{JSON-RPC 信封与方法分发}

A2A 在传输层选择了 JSON-RPC 2.0 而非 REST。所有方法（\texttt{agent/getCard}、\texttt{tasks/send}、\texttt{tasks/get}、\texttt{tasks/cancel}、\texttt{tasks/sendSubscribe}）都通过同一个 \texttt{POST /a2a} 端点提交，靠请求体里的 \texttt{method} 字段路由。代码~\ref{lst:jsonrpc} 给出信封类型。

\begin{lstlisting}[language=Rust, caption={JSON-RPC 信封类型（fuxi-a2a/src/jsonrpc.rs）}, label={lst:jsonrpc}]
pub const JSONRPC_VERSION: &str = "2.0";

pub mod method {
    pub const AGENT_GET_CARD: &str = "agent/getCard";
    pub const TASKS_SEND: &str = "tasks/send";
    pub const TASKS_SEND_SUBSCRIBE: &str = "tasks/sendSubscribe"; // 流式
    pub const TASKS_GET: &str = "tasks/get";
    pub const TASKS_CANCEL: &str = "tasks/cancel";
}

#[derive(Debug, Clone, Serialize, Deserialize)]
pub struct JsonRpcRequest {
    pub jsonrpc: String,
    pub id: serde_json::Value,        // 多态：string/number/null
    pub method: String,
    pub params: Option<serde_json::Value>,
}

#[derive(Debug, Clone, Serialize, Deserialize)]
pub struct JsonRpcResponse {
    pub jsonrpc: String,
    pub id: serde_json::Value,
    pub result: Option<serde_json::Value>,  // 成功
    pub error: Option<JsonRpcError>,        // 失败
}

#[derive(Debug, Clone, Serialize, Deserialize)]
pub struct JsonRpcError {
    pub code: i32,
    pub message: String,
    pub data: Option<serde_json::Value>,
}
\end{lstlisting}

\texttt{params} 与 \texttt{result} 都用 \texttt{serde\_json::Value} 而非强类型——这样 transport 层就和具体业务方法解耦：jsonrpc 模块只关心信封是否合法，业务侧再把 \texttt{Value} 反序列化为 \texttt{SendTaskRequest}/\texttt{Task} 等具体类型。这种分层在协议演进时非常有用：未来 A2A 加新方法或字段，jsonrpc 层无需变更，业务层各自适配即可。错误码遵循 JSON-RPC 标准：\texttt{-32700} 解析失败、\texttt{-32601} 方法未找到、\texttt{-32602} 参数错、\texttt{-32603} 内部错。

JSON-RPC 选择性地放弃了 RESTful 的资源化语义，但换来了三个对 agent 互通极有价值的特性。第一是统一信封格式——所有方法的请求与响应结构完全一致，client 只需写一个泛型 \texttt{call<T>(method, params)}，不用为每个端点单独搭 HTTP 调用；这对未来支持十几种异构 agent 至关重要。第二是 \texttt{id} 字段允许 client 异步并发多个请求并按 \texttt{id} 关联响应，无需依赖 HTTP 层的 keepalive/multiplexing。第三是错误码与错误对象天然内嵌在响应里，HTTP 层始终返 200，避免了"4xx 表示协议错还是业务错"这种 RESTful 长期争议。代价是 JSON-RPC 没有 OpenAPI 那种成熟的 schema 文档生态，但 A2A 规范里 \texttt{TaskState} / \texttt{Part} / \texttt{AgentCard} 等核心类型本身就是 schema 替代物，加上 fuxi 自带的 \texttt{wire.rs} 里 6 个 \texttt{roundtrip} 测试，这一缺陷在工程上是可接受的。

\texttt{id} 字段类型用 \texttt{serde\_json::Value} 而非 \texttt{String}——这是因为 JSON-RPC 2.0 规范允许 id 是 string、number 或 null 三种形态，硬约束成 String 会破坏与 Python/JavaScript client 的互通（这些语言生成的 id 经常是 number）。\texttt{Value} 类型让 server 端按字面值原样回传，无需解析数值或字符串。这种"协议层 id 不感知具体类型"的设计是 JSON-RPC 实现里的标准实践，但凡有 client 用整数 id 而 server 强制 String，就会出现"server 永远返不出 id 让 client 认得出来"的诡异现象。

\subsubsection{Server 端 Axum 路由与 SSE 升级}

服务端基于 Axum 0.7 实现。整个 server 模块的关键设计是"单一端点 + method 分发"——只有一个 \texttt{POST /a2a} 路由，所有方法都到这里报到。代码~\ref{lst:a2a-router} 是 router 装配。

\begin{lstlisting}[language=Rust, caption={A2A server router 装配（fuxi-a2a/src/server.rs）}, label={lst:a2a-router}]
#[async_trait]
pub trait A2AService: Send + Sync + 'static {
    async fn agent_card(&self) -> Result<AgentCard>;
    async fn send_task(&self, req: SendTaskRequest) -> Result<Task>;
    async fn send_task_subscribe(
        &self,
        req: SendTaskRequest,
    ) -> Result<BoxStream<'static, ServerSentEvent>>;
    async fn get_task(&self, id: &str) -> Result<Task>;
    async fn cancel_task(&self, id: &str) -> Result<Task>;
}

pub fn router<S: A2AService>(service: Arc<S>) -> Router {
    Router::new()
        .route("/a2a", post(dispatch::<S>))
        .with_state(service)
}
\end{lstlisting}

\texttt{A2AService} trait 是协议层与业务层的契约——任何 Rust agent 实现想暴露 A2A 端点只需实现这五个方法，HTTP/JSON-RPC/SSE 编解码全部由本 crate 的 server 模块处理。这种"协议库提供 trait 由调用方实现"的反向控制 inversion of control 是标准 axum 服务的扩展模式。

代码~\ref{lst:a2a-dispatch} 是单入口分发函数，体现了"普通方法返 JSON / 订阅方法升级 SSE"的双形态融合。

\begin{lstlisting}[language=Rust, caption={A2A server 单入口分发与 SSE 升级（fuxi-a2a/src/server.rs）}, label={lst:a2a-dispatch}]
async fn dispatch<S: A2AService>(
    State(service): State<Arc<S>>,
    body: axum::body::Bytes,
) -> Response {
    let req: JsonRpcRequest = match serde_json::from_slice(&body) {
        Ok(r) => r,
        Err(e) => return json_err(Value::Null, Error::CODE_PARSE,
                                   format!("parse: {e}"),
                                   StatusCode::BAD_REQUEST),
    };

    match req.method.as_str() {
        method::AGENT_GET_CARD => {
            handle_simple(req.id, service.agent_card().await).await
        }
        method::TASKS_SEND => {
            let params: SendTaskRequest = parse_params(req.params.clone())?;
            handle_send(req.id, service.send_task(params).await).await
        }
        method::TASKS_SEND_SUBSCRIBE => {
            let params: SendTaskRequest = parse_params(req.params.clone())?;
            match service.send_task_subscribe(params).await {
                Ok(s)  => into_sse(s).into_response(),  // 升级 SSE
                Err(e) => jsonrpc_error(req.id, e),
            }
        }
        method::TASKS_GET    => { /* ... */ }
        method::TASKS_CANCEL => { /* ... */ }
        other => json_err(req.id, Error::CODE_METHOD_NOT_FOUND,
                          format!("unknown method: {other}"),
                          StatusCode::OK),
    }
}

fn into_sse(
    s: BoxStream<'static, ServerSentEvent>,
) -> Sse<impl Stream<Item = Result<AxumSse, Infallible>>> {
    let mapped = s.map(|ev| {
        let axum_ev = AxumSse::default()
            .event(ev.event)
            .data(ev.data.to_string());
        Ok::<_, Infallible>(axum_ev)
    });
    Sse::new(mapped).keep_alive(KeepAlive::default())
}
\end{lstlisting}

设计精妙处在于：\texttt{tasks/sendSubscribe} 和其他四个方法共用同一个 \texttt{POST /a2a} URL，但响应类型不同——前者返回 \texttt{text/event-stream}，后者返回 \texttt{application/json}。Axum 的 \texttt{IntoResponse} trait 让两种返回类型在同一个函数签名下统一，调用方只需检查 \texttt{Content-Type} 即可决定如何消费。Google 规范允许两种实现风格："单端点 + method 分发 + SSE 升级"和"双端点 + 显式 \texttt{/a2a/stream}"，伏羲选了前者，因为这能简化 reverse proxy 配置，nginx 只需转发一个 location；\texttt{KeepAlive::default()} 启用心跳防止中间代理（cloudflare 等）以"60 秒空闲"为由切断长连接。

\texttt{into\_sse} 这个 helper 体现了"协议层与框架层解耦"的工程价值。\texttt{A2AService::send\_task\_subscribe} 返回类型是 \texttt{BoxStream<'static, ServerSentEvent>}，其中 \texttt{ServerSentEvent} 是伏羲自定义的纯数据结构（含 \texttt{event: String} 与 \texttt{data: serde\_json::Value} 两个字段），完全不依赖 axum；\texttt{into\_sse} 把这个流翻译成 axum 的 \texttt{Sse<Stream>} 类型再 \texttt{into\_response()}。这种翻译层让业务实现（cc adapter / codex adapter）写完全 framework-agnostic 的代码——它们可以单独单元测试，无需启动 axum server；将来如果换 HTTP 框架（hyper raw / actix-web），只要重写这个 helper 即可。这是一种朴素但极有效的"端口与适配器"模式，比引入完整的 hexagonal architecture 框架要轻得多。

错误处理上，\texttt{handle\_simple} 与 \texttt{handle\_send} 两个 helper 把 \texttt{Result<T>} 转成 \texttt{JsonRpcResponse}：成功路径直接 \texttt{serde\_json::to\_value(v)} 后包成 \texttt{success(id, val)}；失败路径调 \texttt{Error::jsonrpc\_code()} 拿到对应错误码（业务错统一映射到 \texttt{-32603 INTERNAL}）后包成 \texttt{failure(id, err)}。整个 server 模块没有一处 \texttt{unwrap()}——所有 \texttt{Err} 都通过 \texttt{?} 或显式 match 转成 JSON-RPC error，HTTP 层始终返 200（除了最初信封解析失败那种连 \texttt{id} 都拿不到的情况，给 400 让传输层也能感知）。这种"错误统统包进 body 而不抛 HTTP 4xx/5xx"的风格是 JSON-RPC 规范要求的，也方便 client 用统一逻辑处理：永远先解析 body，再看 \texttt{error} 字段。

\subsubsection{Client 端 Reqwest 调用与 SSE 帧解析}

客户端基于 Reqwest 实现。代码~\ref{lst:a2a-client-send} 是 \texttt{send\_task} 的同步调用形态。

\begin{lstlisting}[language=Rust, caption={A2A client send\_task 调用（fuxi-a2a/src/client.rs）}, label={lst:a2a-client-send}]
pub struct A2AClient {
    endpoint: Url,
    http: reqwest::Client,
}

impl A2AClient {
    pub async fn send_task(&self, req: SendTaskRequest) -> Result<Task> {
        let envelope = JsonRpcRequest::new(
            new_id(), method::TASKS_SEND, json!(req));
        let resp = self.http
            .post(self.endpoint.clone())
            .json(&envelope)
            .send().await?;
        let body: JsonRpcResponse = resp.json().await?;
        if let Some(err) = body.error {
            return Err(Error::JsonRpc {
                code: err.code,
                message: err.message,
            });
        }
        let result = body.result
            .ok_or(Error::A2A("missing result".into()))?;
        let resp: SendTaskResponse = serde_json::from_value(result)?;
        Ok(resp.task)
    }
}
\end{lstlisting}

这条调用路径是同步的：发请求、收完整 JSON、返回 \texttt{Task}。但 \texttt{tasks/sendSubscribe} 必须流式接收，因为 long-running agent 可能花几分钟才完成、中途持续推送 \texttt{working} 状态与 \texttt{artifact} 增量。代码~\ref{lst:a2a-sse-parse} 是关键的 SSE 帧解析器。

\begin{lstlisting}[language=Rust, caption={A2A client SSE 帧解析（fuxi-a2a/src/client.rs）}, label={lst:a2a-sse-parse}]
fn parse_sse_stream<S>(byte_stream: S)
    -> impl Stream<Item = Result<ServerSentEventPayload>> + Send
where
    S: Stream<Item = std::result::Result<bytes::Bytes, reqwest::Error>>
        + Send + 'static,
{
    async_stream::stream! {
        let mut bs = Box::pin(byte_stream);
        let mut buf = String::new();
        while let Some(chunk) = bs.next().await {
            let chunk = chunk?;
            buf.push_str(std::str::from_utf8(&chunk)?);
            // 处理已到达的 \n\n 帧——可能一个 chunk 含 0/1/N 个完整帧
            while let Some(idx) = buf.find("\n\n") {
                let frame = buf[..idx].to_string();
                buf.drain(..idx + 2);
                match decode_sse_frame(&frame) {
                    Ok(Some(payload)) => yield Ok(payload),
                    Ok(None) => continue,
                    Err(e)   => yield Err(e),
                }
            }
        }
    }
}

fn decode_sse_frame(frame: &str) -> Result<Option<ServerSentEventPayload>> {
    let mut event_name: Option<String> = None;
    let mut data_lines: Vec<&str> = Vec::new();
    for line in frame.split('\n') {
        if line.is_empty() || line.starts_with(':') { continue; }
        if let Some(rest) = line.strip_prefix("event:") {
            event_name = Some(rest.trim().to_string());
        } else if let Some(rest) = line.strip_prefix("data:") {
            data_lines.push(rest.trim_start_matches(' '));
        }
    }
    let data = data_lines.join("\n");
    match event_name.as_deref().unwrap_or("message") {
        EVENT_STATUS   => Ok(Some(ServerSentEventPayload::Status(
                                  serde_json::from_str(&data)?))),
        EVENT_ARTIFACT => Ok(Some(ServerSentEventPayload::Artifact(
                                  serde_json::from_str(&data)?))),
        other => Err(Error::A2A(format!("unknown event: {other}"))),
    }
}
\end{lstlisting}

实现的关键是"增量缓冲 + 双换行分帧"。\texttt{reqwest::Response::bytes\_stream()} 返回的是 TCP chunk 级流，单条 SSE 消息（\texttt{event: status\textbackslash ndata: \{...\}\textbackslash n\textbackslash n}）可能跨越多个 chunk，也可能一个 chunk 内塞了好几条完整消息。\texttt{buf} 字符串持续累加，每次新 chunk 进来后用 \texttt{find("\textbackslash n\textbackslash n")} 切出所有完整帧，剩余不完整部分留待下次。SSE 规范\cite{whatwg2024sse}还要求兼容注释行（\texttt{:} 开头）与多行 \texttt{data:}（拼接），\texttt{decode\_sse\_frame} 严格按规范处理。这段代码乍看不长，但它是流式 agent 通讯的"最后一公里"，TCP chunk 边界处理稍有疏忽就会出现"消息卡住直到下一条到来"的诡异行为。

性能上，本实现选了 \texttt{String} 增量缓冲 + \texttt{find("\textbackslash n\textbackslash n")} 这种 O(n) 重扫的朴素策略，而非 \texttt{tokio\_util::codec::FramedRead} 这类工业级 codec。原因是 SSE 帧体积通常很小（fuxi 的 status 事件 < 4KB，artifact 事件 < 64KB），单帧重扫开销可忽略；而 codec 路线需要引入 \texttt{tokio\_util} 整套依赖、定义 \texttt{Decoder} trait，代码量翻倍但收益微乎其微。这种「按场景选择实现复杂度」的判断在伏羲实现里反复出现——A2A crate 不引入 hyper raw 而用 reqwest，不引入私有 HTTP 栈而完全依赖 axum，都是同一个工程取舍：协议层定位为产品接口而非性能瓶颈，可维护性与互通性优先于极致性能。

测试覆盖上，\texttt{tests/roundtrip.rs} 包含三组关键测试：\texttt{roundtrip\_all\_methods} 用一个内置 echo agent 跑完 5 个 RPC 方法的完整调用闭环，验证 client 与 server 的对称性；\texttt{roundtrip\_subscribe\_stream} 在 server 端故意把 SSE 帧切成 1 字节 chunks 发出，验证 client 端 \texttt{parse\_sse\_stream} 在最恶劣 chunk 边界下的正确性；\texttt{unknown\_method\_returns\_jsonrpc\_error} 验证 server 对未知方法返回 200 + \texttt{-32601} error 而非 HTTP 4xx——这条 invariant 是 JSON-RPC 规范要求的，违反它会让标准 JSON-RPC client 库困惑。\texttt{wire.rs} 内还有 6 个针对 wire 数据结构的 \texttt{roundtrip} 测试，重点是验证 \texttt{TaskState::InputRequired} 序列化成 \texttt{"input-required"}（而非 \texttt{"InputRequired"}）以及 \texttt{From<core::AgentCard>} 转换的字段映射正确性。三层测试加起来 358 行，相当于核心代码的 30\%——一个协议库应有的测试密度。

\subsubsection{InputRequired 状态：fuxi 对 A2A 规范的扩展}

如前所述，\texttt{InputRequired} 是伏羲对 A2A 1.0 的实质性扩展。其语义是：当 agent 在执行过程中遇到必须由人工决策的分支（典型例子是 reviewer 看完 PR 后不能自己 merge），不发 \texttt{Failed} 也不发 \texttt{Completed}，而是发一条 \texttt{TaskState::InputRequired} 状态事件并暂停推送，等外部编排器（玄女）通过 \texttt{tasks/send} 再发一条用户 message 才能继续。这一设计在 fuxi 内部对应 \texttt{Task::PendingApproval} 状态，在 wire 上映射成 \texttt{"input-required"}，对端编排器看到这个状态可以决策是否把对话权转交给真人用户、是否启动其他 agent 协作。这种"人在回路"显式建模在多 agent 协同系统中不是可有可无的奢侈品，而是工程必要——否则 agent 在等待人工时只能伪装成"超长 working"，与真正的死循环无法区分。

更深一层看，\texttt{InputRequired} 是 fuxi 整套设计哲学的微缩展示。Google 规范默认 agent 是自治闭环系统，要么成功要么失败；伏羲的核心立场是"agent 是工具不是替代品"，必须给人工决策点显式建模。这一立场延伸到伏羲的多个其它设计——玄女永远有知情权（公理 \#2）、deliverable 必须显式 produce 而非自动推断、跨节点任务的 \texttt{InputRequired} 会被 republish 到 home 节点的 EventBus 让玄女感知。\texttt{InputRequired} 既是协议层的扩展，也是产品价值观的表达。这种"协议形态承载产品哲学"的设计模式，在长寿协议（如 SMTP 的 RFC 5321 通过 EHLO 扩展机制让协议跟得上时代）里反复见过，A2A 作为 LLM 时代的新协议大概率也会沿着同样的轨迹演进；fuxi 在这条路径上的小步先行，可以为后续规范修订提供一个可工作的实证参考。

回看整个 \texttt{fuxi-a2a} crate：1\,063 行核心库 + 258 行测试，覆盖 5 个 RPC 方法、6 种 \texttt{TaskState}、3 种 \texttt{Part} 类型；端到端测试 \texttt{roundtrip\_all\_methods} 验证完整的 client-server 调用闭环，\texttt{roundtrip\_subscribe\_stream} 专门验证 SSE 帧在 TCP chunk 边界上的正确切割。这是 \texttt{fuxi} 平台为本地优先 Agent 协作场景定制的一份 A2A 1.0 实现，是本工作的核心贡献之一。与生态内 a2a-rs、a2a-client、a2a-types 等第三方 crate 处在不同定位上：后者以独立 SDK 形态存在，提供通用协议封装；本 crate 则与 \texttt{fuxi-events} 事件总线、玄女—门客编排层、反向 WebSocket sandbox 同源演进，并在协议层扩展了 \texttt{InputRequired} 人工介入语义。从复杂度看，它跨越了 wire format、JSON-RPC envelope、HTTP server、SSE streaming、reverse client 五个层次的协议工程；从平台价值看，它把 fuxi 的可观测性约束与编排层假设直接编码进协议实现。本 crate 的实现质量不仅体现在代码量上，还体现在它经受住了与多个真实 agent endpoint 的互通验证——伏羲的 cc 适配器、codex 适配器以及一个外部 Python 参考实现都通过它做端到端调用，证明本实现的协议合规性达到了规范级别。

\subsection{分布式通讯实现}

伏羲 v2 引入了跨节点能力：用户的 home 机（玄女常驻处）可以把任务分发到 worker 节点（其他物理机或同机不同 user 沙箱）执行。这套分布式层并非重新发明轮子，而是借鉴了 Spanner 的"逻辑时间 + 一致快照"思路\cite{corbett2012spanner}与 ZooKeeper 的"会话 + watcher"心跳模型\cite{hunt2010zookeeper}，但实现上做了大幅简化——伏羲面向的是"个人开发者多机部署"，不是数据中心规模，因此没有引入 Raft\cite{ongaro2014raft} 等共识协议，而是接受"home 节点是单点协调者"的简化模型。这种简化在工程上非常重要：Raft 等协议需要至少三节点 quorum 才能保证可用性，但个人开发者大多只有一台 mac mini + 一台公司笔记本两个节点，启动 Raft 反而会因为 quorum 不足而完全不可用。fuxi 的设计场景明确是"home 是单点真相，worker 是任意数量的从节点"，这与传统数据中心需要的对称容错模型截然不同，因此简化后的架构反而更契合实际部署。

dist controller 是部署在 home 节点的 axum service，监听同一进程内的两套路由：\texttt{/api/*} 走用户 cookie 鉴权（IM 端使用），\texttt{/dist/*} 走 HMAC-SHA256 鉴权（worker 节点使用）。这种"一套代码两套鉴权"通过 axum 的 \texttt{Router::nest + middleware\_fn} 优雅地实现：\texttt{/api/*} 上挂 cookie middleware，\texttt{/dist/*} 上挂 HMAC middleware，handler 内部完全共享业务逻辑。

HMAC 鉴权是分布式安全性的核心。代码~\ref{lst:hmac-verify} 给出验签实现。

\begin{lstlisting}[language=Rust, caption={HMAC-SHA256 验签实现（fuxi-cli/src/dist\_auth.rs）}, label={lst:hmac-verify}]
pub fn verify_request(
    secret: &HmacSecret,
    method: &str,
    path: &str,
    timestamp_ms: u64,
    nonce: &str,
    body: &[u8],
    provided_sig: &str,
    max_skew_ms: u64,
    nonces: &NonceCache,
) -> Result<(), HmacError> {
    let now = now_unix_ms();
    let skew = now.abs_diff(timestamp_ms);
    if skew > max_skew_ms {
        return Err(HmacError::ClockSkew);   // 重放/慢机
    }
    let expected = sign_request(secret, method, path,
                                 timestamp_ms, nonce, body);
    let expected_bytes = expected.as_bytes();
    let provided_bytes = provided_sig.as_bytes();
    // ct_eq 要求等长——长度不等先短路（长度泄漏不敏感）。
    if expected_bytes.len() != provided_bytes.len() {
        return Err(HmacError::BadSignature);
    }
    // 常量时间比较：防 timing oracle 暴露字节级匹配进度
    if expected_bytes.ct_eq(provided_bytes).unwrap_u8() != 1 {
        return Err(HmacError::BadSignature);
    }
    // 签名验过再查 nonce——避免 attacker 用伪签名灌满 nonce cache
    if !nonces.check_and_insert(nonce) {
        return Err(HmacError::ReplayedNonce);
    }
    Ok(())
}
\end{lstlisting}

实现的三个安全要点：第一，时间戳偏差检查放最前——一旦超过 \texttt{max\_skew\_ms}（默认 5 分钟）直接拒绝，无论签名是否合法；这同时防 replay（旧请求）和反向 replay（用未来时间戳预制请求）。第二，签名比对用 \texttt{subtle::ConstantTimeEq::ct\_eq} 的常量时间实现，而非朴素 \texttt{==}——后者会在第一个不匹配字节立即返回，攻击者可以通过测量响应时间逐字节猜出正确签名（timing oracle）。第三，nonce 检查刻意放在签名验证之后，防止攻击者用大量伪造请求灌满 nonce cache 实现拒绝服务（DoS）；只有签名通过的请求才有资格"消耗" nonce 槽位。

签名的 canonical bytes 由 \texttt{method + path + timestamp + nonce + body} 五个字段拼成。这种"全要素签名"的好处是任何字段被篡改都会让签名失败——攻击者拿到一条历史请求即便 timestamp 在窗口内、nonce 没用过，只要改任何一个字节都过不了。\texttt{NonceCache} 内部是一个有限大小的 LRU + 过期时间双重淘汰结构：超过 \texttt{max\_skew\_ms} 时间窗的 nonce 自动失效（因为时间戳检查会先拦截），同时 LRU 淘汰避免无限增长。HMAC secret 通过 \texttt{FUXI\_DIST\_HMAC\_SECRET} 环境变量注入，daemon 启动时若未设置直接拒绝起 \texttt{/dist/*} 路由——这种"必须显式配置"的 fail-safe 比"默认 secret 等待用户改"更难误用，许多生产事故就源于"默认 admin/admin"被忘掉。

worker 节点与 home 之间有四条路径：\texttt{/dist/register}（worker 上线时声明能力），\texttt{/dist/heartbeat}（每 2 秒上报心跳与 inflight 计数），\texttt{/dist/pull}（worker 主动来 home 取活），\texttt{/dist/event}（worker 把执行过程中产生的 \texttt{Event} 回传 home）。这四条路径全部走 HTTP POST + HMAC 鉴权，不引入 gRPC/WebSocket。设计上有意保持 stateless：home 节点不主动推任务给 worker，而是 worker 周期性 pull——这样 worker 在防火墙后也能工作，无需暴露公网端口；只要 worker 能向外打 HTTPS，就能加入分布式集群。inflight 计数与心跳一起上报，home 端的 \texttt{NodeLoadProvider} 用它做负载感知调度：\texttt{saturation = inflight / max\_concurrency}，\texttt{auto\_pin} 时优先选最低 saturation 的节点。

\texttt{auto\_pin\_from\_project} 还有一条隐性 invariant：当用户显式 \texttt{@<node>} pin 时，绝不能悄悄改写。这条规则在 \texttt{Fuxi::auto\_pin\_from\_project} 入口立即检查 \texttt{task.pinned\_node.is\_some()}——只要已 pin 就直接返 None；调用方再加 \texttt{task.pinned\_node.is\_none()} 守卫双保险。这种"用户显式 > 系统智能"的约束在分布式调度中至关重要，否则用户会在"任务为什么没在我指定的机器上跑"的问题上反复 debug 而看不到根因。

worker 节点 inflight 回收靠 \texttt{WorkerHeartbeat} 与 \texttt{WorkerStaleSwept} 两个事件实现。home 端订阅 EventBus 的 watcher 维护一个 \texttt{Map<NodeId, last\_seen>}，超过 30 秒未收到心跳的 worker 被标记为 stale，其上 inflight 的任务全部 republish 为 \texttt{TaskFailed\{ cause: "worker\_stale" \}}，让玄女决定是否重试。这种基于事件的 stale 检测完全异步、不轮询，符合公理 \#3。

dist controller 还需处理一个分布式系统普遍头疼的问题：worker 短暂掉线（VPN 抖动 / SSH 隧道断开 / wifi 切换）期间，home 端是直接判 stale 触发任务回收，还是给一个宽容窗口？伏羲选择了相对宽松的 30 秒——在 2 秒心跳间隔下相当于允许丢 14 次心跳。这个数字是 trade-off 的产物：过短会让用户在地铁上断网 5 秒就丢任务，过长会让真正死掉的 worker 长时间占着 inflight 槽。worker 端在 \texttt{spawn\_worker} 时还有一个 best-effort 的 \texttt{git fetch} 步骤——把上游默认分支拉到本地以减少后续编译时的网络往返，但任何失败只 warn 不 fail，避免短暂网络抖动让任务凭空消失。这种"基础设施操作做 best-effort + 业务操作做 strict"的二分法在分布式实现里反复出现：基础操作能容错就容错，业务操作必须明确成功或明确失败，绝不允许"半成功"暧昧状态。

\subsection{本章小结}

本章把第三章的设计落到了代码。我们看到 fuxi-events 用"非阻塞 publish + lag 哨兵"调和了实时性与持久化的矛盾，看到 fuxi-orchestrator 用 publish-then-pump 模式把 agent 的私有 receiver 转换为全局事件流，看到 fuxi-a2a 从零实现了 Rust 生态尚不存在的 A2A 1.0 完整客户端与服务端、并通过 \texttt{InputRequired} 状态扩展了 Google 规范以支持人在回路场景，看到 dist controller 用 HMAC-SHA256 + nonce cache + 常量时间比较保证跨节点通讯的安全性。所有代码片段均直接节选自 \texttt{crates/} 下的真实实现，可在主仓 HEAD \texttt{04eb34c} 处对位审阅。

总结这一章的三条工程经验：第一，\textbf{协议层与业务层显式分离}——fuxi-a2a 只关心 wire 与 envelope，业务侧通过 \texttt{A2AService} trait 注入；fuxi-events 只关心持久化与扇出，业务事件类型由 fuxi-core 定义。这种分层让协议升级与业务演进互不干扰，每个 crate 都能独立测试与迭代。第二，\textbf{非阻塞 + 哨兵 > 阻塞 + 完美}——实时系统下，让 publisher 阻塞等待写库会拖垮整条调用链；用后台 spawn + lag 哨兵实现"原始 publisher 永远不阻塞、订阅者能感知 backpressure、事件永不丢失"三者兼顾。第三，\textbf{用户显式 \(>\) 系统智能}——\texttt{auto\_pin\_from\_project} 不能改写已 pin 的 task；\texttt{shutdown\_agent} 必须豁免玄女；这些"看起来多余的守卫"是用户信任系统的基础。下一章将给出在这套实现之上做的端到端实测——延迟、吞吐、错误恢复、跨节点协作场景的量化数据，验证设计目标是否达成。
